%=============================================================================!
% 11.4  MULTIPLE TIME STEPS                                                   !
%=============================================================================!
\section{Multiple time steps}
\label{sec:cs-respa}

Constraints remove the fast motions. A second route keeps every motion
and recognises that the forces behind them differ, in how fast they change
and in what they cost. In the flexible water of \cref{sec:cs-water}, the
springs inside each molecule change over the \SI{7.33}{\femto\second}
period of its fastest vibration and cost almost nothing to compute: one
call for all 64 molecules takes \SI{0.04}{\milli\second}. The
Lennard-Jones and Ewald forces between molecules change mainly as whole
molecules move and turn, and one call takes between 23 and
\SI{26}{\milli\second} in our measurements, some 600 times as long. Velocity Verlet
computes both at every step of \SI{0.41}{\femto\second}. This section
builds a method that computes the costly forces less often.

\subsection*{A splitting within a splitting}

\Cref{sec:in-splitting} split the Liouville operator into a part
$\Liou_{\mathrm r}$ that moves the positions and a part $\Liou_{\mathrm
p}$ that moves the momenta by the forces. The forces are a sum, fast plus
slow, $\vb{F}_i = \vb{F}_i^{\mathrm f} + \vb{F}_i^{\mathrm s}$, and since
$\Liou_{\mathrm p}A = \sum_i\vb{F}_i\cdot\grad_{\vb{p}_i}A$ contains each
force to the first power, it splits in the same way, $\Liou_{\mathrm p} =
\Liou_{\mathrm p}^{\mathrm f} + \Liou_{\mathrm p}^{\mathrm s}$, each part a
kick by its own forces alone. The derivation of \cref{eq:in-trotter}
used nothing about $X$ and $Y$ beyond keeping the order of their products,
so with $X = \Liou_{\mathrm p}^{\mathrm s}$ and $Y = \Liou_{\mathrm r} +
\Liou_{\mathrm p}^{\mathrm f}$ it gives
\begin{equation*}
   \mathrm{e}^{\dt\Liou} = \mathrm{e}^{\frac12\dt\Liou_{\mathrm p}^{\mathrm s}}
   \,\mathrm{e}^{\dt(\Liou_{\mathrm r} + \Liou_{\mathrm p}^{\mathrm f})}\,
   \mathrm{e}^{\frac12\dt\Liou_{\mathrm p}^{\mathrm s}} + \order{\dt^3} :
\end{equation*}
half a kick by the slow forces, the motion for $\dt$ under the fast forces
alone, and half a kick by the slow forces. The middle factor is the
motion of a system with only the cheap forces, which $n$ velocity Verlet
steps of $\dt/n$ follow,
\begin{equation*}
   \mathrm{e}^{\dt(\Liou_{\mathrm r} + \Liou_{\mathrm p}^{\mathrm f})}
   = \Bigl(\mathrm{e}^{\frac{\dt}{2n}\Liou_{\mathrm p}^{\mathrm f}}\,
   \mathrm{e}^{\frac{\dt}{n}\Liou_{\mathrm r}}\,
   \mathrm{e}^{\frac{\dt}{2n}\Liou_{\mathrm p}^{\mathrm f}}\Bigr)^n
   + \order{\dt^3/n^2} ,
\end{equation*}
the error of $n$ steps each in error by order $(\dt/n)^3$, $n(\dt/n)^3
= \dt^3/n^2$ in all. The outer
step is $\dt$, taken with one call of the slow forces; the inner step is
$\dt/n$, taken $n$ times with one call of the fast forces each. Every factor is a
kick or a drift, so the product, like velocity Verlet, keeps the full
symplectic condition (\cref{sec:in-reversal}), and since it reads the same
from either end, reversing the velocities and stepping again undoes it
factor by factor, as for velocity Verlet in \cref{sec:in-reversal}: it is
reversible. Tuckerman, Berne and Martyna derived the
method in this way, as the reversible reference system propagator
algorithm, RESPA\cite{Tuckerman1992}. With $n = 1$ it is velocity Verlet
with the two forces added (\cref{ex:cs-respa-one}).

The function \texttt{respa\_step} of \texttt{mdlab.respa} is the outer
step:
\begin{code}
    v += 0.5 * dt * force_to_accel * np.asarray(slow_forces) / m
    u_fast = 0.0
    for _ in range(n_inner):
        v += 0.5 * h * force_to_accel * f_fast / m
        r += h * v
        u_fast, f_fast = fast(r)
        v += 0.5 * h * force_to_accel * f_fast / m
    u_slow, f_slow = slow(r)
    v += 0.5 * dt * force_to_accel * f_slow / m
\end{code}
\noindent where \texttt{h} is the inner step $\dt/n$ and \texttt{f\_fast}
holds the fast forces at the start of the step, passed in from the step
before. Its test checks
that with one inner step it gives the
positions and velocities of velocity Verlet to $10^{-12}$. For water,
\texttt{WaterModel} supplies the two parts: \texttt{intramolecular}, the
springs, as the fast forces, and \texttt{intermolecular}, Lennard-Jones
and Ewald, as the slow ones.

\subsection*{What it buys}

\Cref{fig:cs-respa} runs the flexible box with inner steps of
\SI{0.25}{\femto\second}, at which velocity Verlet alone keeps the energy
fluctuation at 0.36\% of the kinetic energy's, and outer steps from
\SI{0.25}{\femto\second} to \SI{4.5}{\femto\second}. At an outer step of
\SI{1}{\femto\second} the fluctuation is 1.8\%, where velocity Verlet at
the same step gives 6.7\%. At the 1\% level RESPA reaches an outer step
of \SI{0.72}{\femto\second} against velocity Verlet's
\SI{0.41}{\femto\second}, 1380 calls of the slow forces per picosecond
against 2430 (from the steps before rounding), a factor of 1.76; at 2\% the factor is 1.82
(\cref{fig:cs-respa}b). Since the slow forces take nearly all the time,
the run is faster by close to that factor.

\begin{figure}[htb]
\centering
\includegraphics{ch11_constraints/respa}
\caption{Multiple time steps for 64 molecules of flexible water: the
springs are the fast forces, followed with inner steps of
\SI{0.25}{\femto\second}; Lennard-Jones and Ewald forces are the slow
ones, applied once per outer step $\dt$. (a) The standard deviation of the
total energy over that of the kinetic energy, against $\dt$, for RESPA
and for velocity Verlet taking all forces at every step $\dt$ (dashed);
dotted, a quarter, a third and a half of the period $\mathcal{T} =
\SI{7.33}{\femto\second}$ of the fastest vibration, and half the next,
$\mathcal{T}' = \SI{8.88}{\femto\second}$. (b) The same fluctuation
against the number of calls of the slow forces per picosecond, which sets
the cost of a run.}
\label{fig:cs-respa}
\end{figure}

The gain stops early. Beyond an outer step of about
\SI{1.8}{\femto\second}, a quarter of the fastest period, the fluctuation
climbs much faster than $\dt^2$, to 8.8\% at \SI{2}{\femto\second} and
43\% at \SI{3}{\femto\second}, and the run blows up at
\SI{4.5}{\femto\second}. Following the fast motion more closely does not
help: inner steps of \SI{0.125}{\femto\second} with an outer step of
\SI{2}{\femto\second} give 10.4\%, no better than 8.8\%. What fails is
the outer step itself, and \cref{sec:cs-resonance} finds why.

\begin{keyidea}
RESPA splits the forces into fast and slow parts and the Liouville
operator with them: half a kick by the slow forces, $n$ velocity Verlet
steps of $\dt/n$ under the fast forces alone, and half a kick by the slow
forces. It is reversible and symplectic. On flexible water it cuts the
calls of the costly forces by a factor of about 1.8 at equal accuracy, and
fails beyond outer steps of about a quarter of the fastest period.
\end{keyidea}

\notebook{11}{split the forces of water and choose the outer step}

\begin{selfcheck}
With \SI{24}{\milli\second} and \SI{0.04}{\milli\second} per call, about how long does one picosecond of
flexible water take by velocity Verlet at \SI{0.41}{\femto\second}, and by
RESPA at an outer step of \SI{0.72}{\femto\second} with three inner
steps?
\end{selfcheck}
